\documentclass{amsart}
% For dealing with references we use the comment environment
\usepackage{verbatim}
\newenvironment{reference}{\comment}{\endcomment}
%\newenvironment{reference}{}{}
\newenvironment{slogan}{\comment}{\endcomment}
\newenvironment{history}{\comment}{\endcomment}

% For commutative diagrams we use Xy-pic
\usepackage[all]{xy}

% We use 2cell for 2-commutative diagrams.
\xyoption{2cell}
\UseAllTwocells

% We use multicol for the list of chapters between chapters
\usepackage{multicol}

% This is generally recommended for better output
\usepackage{lmodern}
\usepackage[T1]{fontenc}

% For cross-file-references

% Package for hypertext links:
\usepackage{hyperref}

% For any local file, say "hello.tex" you want to link to please

% Theorem environments.
%
\theoremstyle{plain}
\newtheorem{theorem}[subsection]{Theorem}
\newtheorem{proposition}[subsection]{Proposition}
\newtheorem{lemma}[subsection]{Lemma}

\theoremstyle{definition}
\newtheorem{definition}[subsection]{Definition}
\newtheorem{example}[subsection]{Example}
\newtheorem{exercise}[subsection]{Exercise}
\newtheorem{situation}[subsection]{Situation}

\theoremstyle{remark}
\newtheorem{remark}[subsection]{Remark}
\newtheorem{remarks}[subsection]{Remarks}

\numberwithin{equation}{subsection}

% Macros
%
\def\lim{\mathop{\mathrm{lim}}\nolimits}
\def\colim{\mathop{\mathrm{colim}}\nolimits}
\def\Spec{\mathop{\mathrm{Spec}}}
\def\Hom{\mathop{\mathrm{Hom}}\nolimits}
\def\Ext{\mathop{\mathrm{Ext}}\nolimits}
\def\SheafHom{\mathop{\mathcal{H}\!\mathit{om}}\nolimits}
\def\SheafExt{\mathop{\mathcal{E}\!\mathit{xt}}\nolimits}
\def\Sch{\mathit{Sch}}
\def\Mor{\mathop{\mathrm{Mor}}\nolimits}
\def\Ob{\mathop{\mathrm{Ob}}\nolimits}
\def\Sh{\mathop{\mathit{Sh}}\nolimits}
\def\NL{\mathop{N\!L}\nolimits}
\def\CH{\mathop{\mathrm{CH}}\nolimits}
\def\proetale{{pro\text{-}\acute{e}tale}}
\def\etale{{\acute{e}tale}}
\def\QCoh{\mathit{QCoh}}
\def\Ker{\mathop{\mathrm{Ker}}}
\def\Im{\mathop{\mathrm{Im}}}
\def\Coker{\mathop{\mathrm{Coker}}}
\def\Coim{\mathop{\mathrm{Coim}}}

% Boxtimes
%
\DeclareMathSymbol{\boxtimes}{\mathbin}{AMSa}{"02}

%
% Macros for moduli stacks/spaces
%
\def\QCohstack{\mathcal{QC}\!\mathit{oh}}
\def\Cohstack{\mathcal{C}\!\mathit{oh}}
\def\Spacesstack{\mathcal{S}\!\mathit{paces}}
\def\Quotfunctor{\mathrm{Quot}}
\def\Hilbfunctor{\mathrm{Hilb}}
\def\Curvesstack{\mathcal{C}\!\mathit{urves}}
\def\Polarizedstack{\mathcal{P}\!\mathit{olarized}}
\def\Complexesstack{\mathcal{C}\!\mathit{omplexes}}
% \Pic is the operator that assigns to X its picard group, usage \Pic(X)
% \Picardstack_{X/B} denotes the Picard stack of X over B
% \Picardfunctor_{X/B} denotes the Picard functor of X over B
\def\Pic{\mathop{\mathrm{Pic}}\nolimits}
\def\Picardstack{\mathcal{P}\!\mathit{ic}}
\def\Picardfunctor{\mathrm{Pic}}
\def\Deformationcategory{\mathcal{D}\!\mathit{ef}}

\setlength{\emergencystretch}{3em}
\begin{document}
\title[Stacks Project, Tag 0D6T]{658 --- Eventually constant cohomology survives derived inverse limits}
\maketitle
\noindent Copyright (C) 2005--2025 Johan de Jong. Stacks Project authors. Source: \href{https://stacks.math.columbia.edu/tag/0D6T}{Tag 0D6T}.

\noindent Editorial difficulty note (not author text). Difficulty H2: The author bounds the spectral-sequence range locally, establishes eventual stabilization of two neighboring hypercohomology groups, and proves injectivity and local surjectivity of the limiting comparison. This local cohomological-dimension argument controls an unbounded limit phenomenon beyond qualification material..

\noindent Editorial provenance note (not author text). History: excerpt from revision \texttt{a0444\allowbreak 6e57e\allowbreak c1fbc\allowbreak 252a8\allowbreak 71afc\allowbreak ec775\allowbreak 2fb28\allowbreak 07b14}, sites-cohomology.tex, lines 5725--5793. Reference presentation was adapted; original-author context and core support blocks were retained or added. The mathematical wording is unchanged. Licensed under GNU FDL 1.2 or later; no invariant sections or cover texts. See ../licenses/COPYING. Context blocks reproduce author definitions and supporting results; they are not additional counted problems.

\begin{situation}
\label{situation-olsson-laszlo}
Let $\mathcal{C}$ be a site. Let $\mathcal{O}$ be a sheaf of rings on
$\mathcal{C}$. Let $\mathcal{A} \subset \textit{Mod}(\mathcal{O})$
be a weak Serre subcategory. We assume the following is true:
there exists a subset $\mathcal{B} \subset \Ob(\mathcal{C})$ such that
\begin{enumerate}
\item every object of $\mathcal{C}$ has a covering whose
members are in $\mathcal{B}$, and
\item for every $V \in \mathcal{B}$ there exists an integer $d_V$
and a cofinal system $\text{Cov}_V$ of coverings of $V$ such
that
$$
H^p(V_i, \mathcal{F}) = 0 \text{ for }
\{V_i \to V\} \in \text{Cov}_V,\ p > d_V, \text{ and }
\mathcal{F} \in \Ob(\mathcal{A})
$$
\end{enumerate}
\end{situation}

\begin{lemma}
\label{lemma-olsson-laszlo-modified}
In Situation \ref{situation-olsson-laszlo} let
$(K_n)$ be an inverse system in $D_\mathcal{A}^+(\mathcal{O})$.
Assume that for every $j$ the inverse system $(H^j(K_n))$
in $\mathcal{A}$ is eventually constant with value $\mathcal{H}^j$. Then
$H^j(R\lim K_n) = \mathcal{H}^j$ for all $j$.
\end{lemma}

\begin{proof}
Let $V \in \mathcal{B}$. Let $\{V_i \to V\}$ be in the set
$\text{Cov}_V$ of Situation \ref{situation-olsson-laszlo}.
Because $K_n$ is bounded below there is a spectral sequence
$$
E_2^{p, q} = H^p(V_i, H^q(K_n))
$$
converging to $H^{p + q}(V_i, K_n)$. See
Derived Categories, Lemma \href{https://stacks.math.columbia.edu/tag/015J}{lemma two ss complex functor (Tag 015J)}.
Observe that $E_2^{p, q} = 0$
for $p > d_V$ by assumption. Pick $n_0$ such that
$$
\begin{matrix}
\mathcal{H}^{j + 1} & = & H^{j + 1}(K_n), \\
\mathcal{H}^j & = & H^j(K_n), \\
\ldots, \\
\mathcal{H}^{j - d_V - 2} & = & H^{j - d_V - 2}(K_n)
\end{matrix}
$$
for all $n \geq n_0$. Comparing the spectral sequences above
for $K_n$ and $K_{n_0}$, we see that for $n \geq n_0$ the
cohomology groups $H^{j - 1}(V_i, K_n)$ and $H^j(V_i, K_n)$
are independent of $n$. It follows that the map on sections
$H^j(R\lim K_n)(V) \to H^j(K_n)(V)$ is injective for $n$ large
enough (depending on $V$), see
Lemma \href{https://stacks.math.columbia.edu/tag/0D6L}{lemma cohomology derived limit injective (Tag 0D6L)}.
Since every object of $\mathcal{C}$ can be covered by elements
of $\mathcal{B}$, we conclude that the map
$H^j(R\lim K_n) \to \mathcal{H}^j$ is injective.

\medskip\noindent
Surjectivity is shown in a similar manner. Namely, pick
$U \in \Ob(\mathcal{C})$ and $\gamma \in \mathcal{H}^j(U)$.
We want to lift $\gamma$ to a section of $H^j(R\lim K_n)$
after replacing $U$ by the members of a covering. Hence we may
assume $U = V \in \mathcal{B}$ by property (1) of
Situation \ref{situation-olsson-laszlo}.
Pick $n_0$ such that
$$
\begin{matrix}
\mathcal{H}^{j + 1} & = & H^{j + 1}(K_n), \\
\mathcal{H}^j & = & H^j(K_n), \\
\ldots, \\
\mathcal{H}^{j - d_V - 2} & = & H^{j - d_V - 2}(K_n)
\end{matrix}
$$
for all $n \geq n_0$. Choose an element $\{V_i \to V\}$ of
$\text{Cov}_V$ such that
$\gamma|_{V_i} \in \mathcal{H}^j(V_i) = H^j(K_{n_0})(V_i)$
lifts to an element $\gamma_{n_0, i} \in H^j(V_i, K_{n_0})$.
This is possible because $H^j(K_{n_0})$ is the sheafification
of $U \mapsto H^j(U, K_{n_0})$ by Lemma \href{https://stacks.math.columbia.edu/tag/0BKV}{lemma sheafification cohomology (Tag 0BKV)}.
By the discussion in the first paragraph of the proof we have that
$H^{j - 1}(V_i, K_n)$ and $H^j(V_i, K_n)$
are independent of $n \geq n_0$. Hence
$\gamma_{n_0, i}$ lifts to an element
$\gamma_i \in H^j(V_i, R\lim K_n)$ by
Lemma \href{https://stacks.math.columbia.edu/tag/0D6K}{lemma RGamma commutes with Rlim (Tag 0D6K)}.
This finishes the proof.
\end{proof}
\end{document}
